\chapter{格点作用量的最简单形式}

\section{规范作用量}

规范作用量可以用闭合圈来表示。我将以阿贝尔 $U(1)$ 模型这一更简单的情形为例来概述其构造，在该模型中链变量是相互对易的复数，而不是 SU(3) 矩阵。考虑最简单的威尔逊圈，即 $1\times1$ 方格 $W_{\mu\nu}^{1\times1}$：
\begin{equation}
\label{eq:plaqexpA}
\begin{aligned}
W_{\mu\nu}^{1\times1} &= U_\mu(x)\, U_\nu(x+\hat{\mu})\, U_\mu^\dagger(x+\hat{\nu})\, U_\nu^\dagger(x) \\
&= e^{iag\big[ A_\mu(x+\hat{\mu}/2) + A_\nu(x+\hat{\mu}+\hat{\nu}/2) -
          A_\mu(x+\hat{\nu}+\hat{\mu}/2) - A_\nu(x+\hat{\nu}/2) \big]}
\end{aligned}
\end{equation}
在 $x + (\hat{\mu}+\hat{\nu})/2$ 附近展开得
\begin{equation}
\label{eq:plaqexpB}
\begin{aligned}
&\exp\Big[ ia^2g(\partial_\mu A_\nu - \partial_\nu A_\mu) +
       \frac{ia^4g}{12}(\partial_\mu^3 A_\nu - \partial_\nu^3 A_\mu) + \ldots \Big] \\
&= 1 + ia^2gF_{\mu\nu} - \frac{a^4g^2}{2}F_{\mu\nu}F^{\mu\nu} + O(a^6) + \ldots
\end{aligned}
\end{equation}
该圈的实部和虚部给出
\begin{equation}
\label{eq:plaqexpC}
\begin{aligned}
\operatorname{Re}\operatorname{Tr}(1-W_{\mu\nu}^{1\times1}) &=
   \frac{a^4g^2}{2}F_{\mu\nu}F^{\mu\nu} + \text{对 } a \text{ 的高阶项} \\
\operatorname{Im}(W_{\mu\nu}^{1\times1}) &= a^2gF_{\mu\nu} .
\end{aligned}
\end{equation}
到目前为止，指标 $\mu$ 和 $\nu$ 尚未收缩。与每个格点相关联的有 $6$ 个正取向互不相同的方格，即 $\{\mu<\nu\}$。对 $\mu,\nu$ 求和，并用额外的 $\frac{1}{2}$ 因子处理双重计数，得到
\[
\frac{1}{g^2}\sum_x\sum_{\mu<\nu}\operatorname{Re}\operatorname{Tr}(1-W_{\mu\nu}^{1\times1})
= \frac{a^4}{4}\sum_x\sum_{\mu,\nu}F_{\mu\nu}F^{\mu\nu}
\to \frac{1}{4}\int d^4x\, F_{\mu\nu}F^{\mu\nu} .
\]
因此，在 $a$ 的最低阶，方格的展开给出连续作用量。非阿贝尔群的推导步骤完全相同，最终表达式（除数值因子外）也一样。以方格定义的 SU(3) 规范作用量的结果称为威尔逊作用量，为
\begin{equation}
\label{eq:Wgaugeaction}
S_g = \frac{6}{g^2}\sum_x\sum_{\mu<\nu}\operatorname{Re}\operatorname{Tr}\,
      \frac{1}{3}(1-W_{\mu\nu}^{1\times1}) .
\end{equation}
由于历史原因，格点计算大多以耦合常数 $\beta \equiv 6/g^2$ 来表示。由于 $\beta$ 也用于许多其他量，例如 $1/kT$ 或 $\beta$-函数，若有可能引起混淆，我将尽量在记号的使用上（$g$ 与 $\beta$）更加明确。

\medskip
\noindent\textit{练习：对物理情形 SU(3) 重复上述泰勒展开。证明
$F_{\mu\nu} = \partial_\mu A_\nu - \partial_\nu A_\mu + g[A_\mu, A_\nu]$ 中的
非阿贝尔项源于 $\lambda$ 矩阵的非对易性质。SU(2) 情形的推导细节见
文献~[8] 的第 3 节。}
\medskip

基于上述格点作用量的构造，有四个要点值得注意。

\textbf{1)} 领头修正为 $O(a^2)$。项
$\frac{a^2}{6}F_{\mu\nu}(\partial_\mu^3 A_\nu - \partial_\nu^3 A_\mu)$
出现在所有平面威尔逊圈的展开中。因此，在经典层面，可以通过选择 $1\times1$
与 $1\times2$ 威尔逊圈（例如）的线性组合作为作用量来消除它，其相对强度由泰勒展开给出（见第 12 节）。

\textbf{2)} 量子效应会产生修正，即 $a^2 \to X(g^2)a^2$，其中在微扰论中
$X(g^2) = 1 + c_1 g^2 + \ldots$，并且会带来额外的非平面圈。因此，作用量的改进将需要包含这些额外的圈，并调整其相对强度，这些强度成为 $g^2$ 的函数。这一点也将在第 12 节中讨论。

\textbf{3)} 用小的圈来定义作用量的原因是计算速度以及减小离散化误差的幅度。例如，$1\times1$ 圈的领头修正正比于 $a^2/6$，而对于 $1\times2$ 圈，它增大到 $5a^2/12$。此外，模拟的成本会增加 $2-3$ 倍。

\textbf{4)} 电场和磁场 $E$ 与 $B$ 正比于 $F_{\mu\nu}$。式~\ref{eq:plaqexpC} 表明它们由威尔逊圈的虚部给出。

\section{费米子作用量}

为了离散化 Dirac 作用量，Wilson 用对称化差分替代了导数，并加入了适当的规范链以保持规范不变性
\begin{equation}
\label{eq:Dslnaive}
\bar\psi\not{D}\psi = \frac{1}{2a}\bar\psi(x)\sum_\mu\gamma_\mu
\big[ U_\mu(x)\psi(x+\hat{\mu}) - U_\mu^\dagger(x-\hat{\mu})
      \psi(x-\hat{\mu}) \big] .
\end{equation}
容易看出，将 $U_\mu$ 和 $\psi(x+\hat{\mu})$ 按格距 $a$ 的幂次做泰勒展开，在 $a \to 0$ 极限下可以恢复 Dirac 作用量。仅保留 $a$ 的领头项，式~\ref{eq:Dslnaive} 变为
\[
\begin{aligned}
&\frac{1}{2a}\bar\psi(x)\gamma_\mu\Big[
   \big(1+iagA_\mu(x+\hat{\mu}/2)+\ldots\big)
   \big(\psi(x)+a\psi'(x)+\ldots\big) \\
&\quad - \big(1-iagA_\mu(x-\hat{\mu}/2)+\ldots\big)
   \big(\psi(x)-a\psi'(x)+\ldots\big) \Big] \\
&= \bar\psi(x)\gamma_\mu\Big(\partial_\mu + \frac{a^2}{6}\partial_\mu^3 +
   \ldots\Big)\psi(x) \\
&\quad + ig\bar\psi(x)\gamma_\mu\Big[A_\mu + \frac{a^2}{2}\Big(
   \frac{1}{4}\partial_\mu^2 A_\mu + (\partial_\mu A_\mu)\partial_\mu +
   A_\mu\partial_\mu^2\Big) + \ldots \Big]\psi(x) ,
\end{aligned}
\]
在 $O(a^2)$ 精度下，这正是欧几里得时空中的标准连续 Dirac 作用量的动能部分。于是便得到最简单的（称为「朴素」）费米子格点作用量
\begin{equation}
\label{eq:naiveact}
\begin{aligned}
\mathcal{S}^N &= m_q\sum_x\bar\psi(x)\psi(x) \\
&\quad + \frac{1}{2a}\sum_x\bar\psi(x)\gamma_\mu\big[
   U_\mu(x)\psi(x+\hat{\mu}) - U_\mu^\dagger(x-\hat{\mu})
   \psi(x-\hat{\mu}) \big] \\
&\equiv \sum_x\bar\psi(x)M_{xy}^N[U]\psi(y)
\end{aligned}
\end{equation}
其中相互作用矩阵 $M^N$ 为
\begin{equation}
\label{eq:naiveactM}
M_{i,j}^N[U] = m_q\delta_{ij} + \frac{1}{2a}\sum_\mu
\big[ \gamma_\mu U_{i,\mu}\delta_{i,j-\mu} -
      \gamma_\mu U_{i-\mu,\mu}^\dagger\delta_{i,j+\mu} \big]
\end{equation}
欧几里得 $\gamma$ 矩阵是厄米的，$\gamma_\mu = \gamma_\mu^\dagger$，且满足 $\{\gamma_\mu,\gamma_\nu\} = 2\delta_{\mu\nu}$。我将采用的表示是
\begin{equation}
\label{eq:gammamatdef}
\vec\gamma = \begin{pmatrix} 0 & i\vec\sigma \\ -i\vec\sigma & 0 \end{pmatrix},
\qquad
\gamma_4 = \begin{pmatrix} 1 & 0 \\ 0 & -1 \end{pmatrix},
\qquad
\gamma_5 = \begin{pmatrix} 0 & 1 \\ 1 & 0 \end{pmatrix} ,
\end{equation}
它与 Bjorken-Drell 约定有如下关系：
$\gamma_i = i\gamma_{BD}^i$，$\gamma_4 = \gamma_{BD}^0$，
$\gamma_5 = \gamma_{BD}^5$。在这一表示中，$\gamma_1, \gamma_3$ 为纯虚数，而 $\gamma_2, \gamma_4, \gamma_5$ 为实数。

泰勒展开表明离散化误差从 $O(a^2)$ 开始。作为另一个简单例证，考虑自由场传播子的逆 $m + \frac{i}{a}\sum_\mu\gamma_\mu\sin(p_\mu a)$。令 $\vec{p}=0$ 并旋转到闵可夫斯基空间（$p_4 \to iE$，即 $\sin p_4 a \to i\sinh Ea$）。然后，利用前向传播子（$\gamma_4$ 的上两个分量），得到
\begin{equation}
\label{eq:Npolemass}
m_q^{\mathrm{pole}} a = \sinh Ea
\end{equation}
这是极点质量与能量之间的关系。这表明，即使在自由场情形下，连续关系 $E(\vec{p}=0) = m$ 也会被 $O(a^2)$ 的修正所破坏。

\noindent\textbf{对称性}：费米子作用量在时空旋转下的不变群是超立方群。完整的欧几里得不变性只在连续极限下才能恢复。作用量在平移 $a$ 以及 $\mathcal{P}$、$\mathcal{C}$ 和 $\mathcal{T}$ 下不变，这可以用表 2 来检验。

朴素作用量 $\bar\psi_x M_{xy}^N\psi_y$ 具有以下整体对称性：
\begin{equation}
\label{eq:vector_psi}
\begin{aligned}
\psi(x) &\to e^{i\theta}\psi(x) \\
\bar\psi(x) &\to \bar\psi(x)e^{-i\theta}
\end{aligned}
\end{equation}
其中 $\theta$ 是连续参数。该对称性与重子数守恒相关，并导致一个守恒的矢量流。对于 $m_q = 0$，作用量在下列变换下也不变
\begin{equation}
\label{eq:axial_psi}
\begin{aligned}
\psi(x) &\to e^{i\theta\gamma_5}\psi(x) \\
\bar\psi(x) &\to \bar\psi(x)e^{i\theta\gamma_5}
\end{aligned}
\end{equation}
在硬截断正则化中同时具有矢量与轴矢对称性将意味着违反 Adler-Bell-Jackiw 定理。事实证明，朴素费米子作用量虽然保持手征对称性，但也具有臭名昭著的费米子加倍问题，如第 9 节所述。这些额外费米子的手征荷恰好抵消了反常。事实上，对 $\mathcal{S}^N$ 的分析导致了 Nielsen-Ninomiya [14] 的 no-go 定理，该定理指出，不可能定义一个局域的、平移不变的、厄米的格点作用量，既保持手征对称性又没有加倍费米子。

我们将讨论对加倍/反常/手征对称性问题的两种解决方案。第一种是在作用量中加入一个额外的项来破坏手征对称性并消除加倍费米子（Wilson 方案）。第二种是利用朴素费米子作用量具有大得多的对称群 $U_V(4)\otimes U_A(4)$ 这一事实，在保持残余手征对称性的同时，将加倍问题从 $2^d = 16 \to 16/4$（交错费米子）。在这两种情况下，都能在连续极限下恢复正确的反常。在深入讨论这些方案之前，我想先讨论积分测度与规范固定的问题，并对 Symanzik 的改进纲领做一概述。

\section{Haar 测度}

通过路径积分将格点QCD表述为量子场论所需的第四个要素，是定义对规范自由度的积分测度。注意，与连续场 $A_\mu$ 不同，格点场是 SU(3) 矩阵，其元素取值有界于 $[0,1]$ 范围内。因此，Wilson 提出用不变群测度，即 Haar 测度，来进行这一积分。该测度的定义为：对群的任意元素 $V$ 和 $W$
\begin{equation}
\label{eq:haarmeasure}
\int dU\, f(U) = \int dU\, f(UV) = \int dU\, f(WU)
\end{equation}
其中 $f(U)$ 是群上的任意函数。这一构造很简单，并且还有一个额外优点：在非微扰研究中，它避免了必须在路径积分中包含规范固定项的问题。这是因为场变量是紧致的。因此不存在发散，可以通过定义下式来归一化测度
\begin{equation}
\label{eq:haarnorm}
\int dU = 1 .
\end{equation}
对于规范不变的量，每个规范拷贝的贡献相同。因此，缺少规范固定项只会给出一个整体的额外数值因子。由于路径积分的归一化，这个因子在关联函数的计算中被抵消。注意，在格点微扰论中必须引入规范固定和 Fadeev-Papov 构造 [36]。

测度和作用量在 $\mathcal{C}$、$\mathcal{P}$、$\mathcal{T}$ 变换下不变，因此，当通过路径积分计算关联函数时，我们必须对组态 $\mathcal{U}$ 以及通过施加 $\mathcal{C}$、$\mathcal{P}$、$\mathcal{T}$ 由它得到的组态求平均。在大多数情况下，这种平均非常自然地发生，例如，欧几里得平移不变性保证了 $\mathcal{P}$ 和 $\mathcal{T}$ 不变性。类似地，取威尔逊圈迹的实部等价于对 $U$ 及其电荷共轭组态求和。关于在由格点离散对称性相联系的组态上求平均后关联函数行为的计算，第 17 节给出了一个可直接套用的方法。

这些对称性的一个重要结果是，只有欧几里得关联函数的实部具有非零的期望值。这对用格点QCD可以计算的量施加了限制。例如，格点QCD 模拟在散射振幅的计算方面并不很成功，因为散射振幅一般是复的 [48]。闵可夫斯基与欧几里得关联函数之间的联系已在第 5 节讨论过。要得到散射振幅，需要先计算所有欧几里得时间 $\tau$ 值下的欧几里得振幅，然后对 $\tau$ 做傅里叶变换，最后维克旋转到闵可夫斯基空间。这样的程序要能成立，需要具有极高精度的数据。遗憾的是，模拟给出的是一组离散 $\tau$ 值上的欧几里得关联函数，且其误差在旋转中被指数放大。目前做到的最好结果是利用 L\"uscher 的方法，该方法把散射相移与有限体积中两粒子能量态的移动联系起来 [27,28]。到目前为止，该方法已在 2--4 维的模型上得到检验，并已用于计算 QCD 中的散射长度 [29,30]。

\section{规范固定与 Gribov 拷贝}

规范不变性与群积分的性质 $\int dU\, U = 0$ 足以证明，只有规范不变的关联函数才具有非零的期望值。这就是著名的 Elitzur 定理 [37]。它指出，定域连续对称性不能被自发破缺。

在某些应用中（夸克与胶子传播子的研究 [38]、三胶子耦合 [39]、强子波函数 [40]，以及利用夸克态确定重整化常数 [41]），在固定的规范下工作是有用的。常见的规范选择有轴规范（$\eta_\mu A_\mu = 0$，其中 $\eta$ 在闵可夫斯基空间中是一个固定的类时矢量，常见选择是 $A_4 = 0$）、库仑规范（$\partial_i A_i = 0$）和朗道规范（$\partial_\mu A_\mu = 0$）。在周期性边界条件下，它们在格点上的对应形式如下。

\textbf{最大轴规范}：在具有周期性边界条件的格点上，不能通过把时间方向的所有链设为 1 来强制 $A_4 = 0$，因为每个方向的威尔逊线都是一个不变量。除了一个格点外，其余格点的规范自由度可以按如下方式消除。在除一个时间片外的所有时间片上，设所有类时链 $U_4 = 1$。在那个时间片上，所有 $U_4$ 链被约束为 4 方向的原始威尔逊线；不过，可以设在除一个 z-平面外的所有 $\hat{z}$ 方向链 $U_3 = 1$。在那个 z-平面上，把 $U_3$ 链设为该方向的威尔逊线，并设在除一条 y-线外的所有 $\hat{y}$ 方向链 $U_2 = 1$。在这条 y-线上，把所有 $U_2$ 链设为 y-线的值，并设在除一个点外的所有 $\hat{x}$ 方向链 $U_1 = 1$。虽然这样的构造是唯一的，但时间片、z-平面、y-线和 x-点的选择以及 $x,y,z,t$ 轴的标记并不是唯一的。此外，这种规范固定不保持欧几里得平移不变性，而且不是光滑的。

\textbf{库仑规范}：格点条件由使下列函数最大化的规范变换 $V(x)$ 给出
\begin{equation}
\label{eq:Coulombgauge}
F[V] = \sum_x\sum_i \operatorname{Re}\operatorname{Tr} V(x)
\bigg( U_i(x)V^\dagger(x+\hat{i}) +
       U_i^\dagger(x-\hat{i})V^\dagger(x-\hat{i}) \bigg)
\end{equation}
该条件在每个时间片上分别成立，且 $i$ 遍历空间指标 $1-3$。

\textbf{朗道规范}：它由与式~\ref{eq:Coulombgauge} 相同的函数定义，只是求和在整个格点上进行，且 $i = 1-4$。

最大化条件，即式~\ref{eq:Coulombgauge}，在连续情形下对应于求 $F[V] = \lVert A^V \rVert = \int d^4x\, \operatorname{Tr}(A_\mu^V)^2(x)$ 的驻点。我们想要的是找到对应于「最光滑」场 $A_\mu$ 的整体极小值；若它不唯一，则找到相应的拷贝。对任何规范组态，通过施加变换集 $\{V(x)\}$ 可以得到无穷多个规范等价组态。这个集合称为规范轨道。第一个问题是，规范固定程序是否给出唯一解，即规范固定算法是否与规范轨道上的起点无关而收敛到同一个最终点。对于最大轴规范，构造是唯一的（对任意 $\{V\}$，组态 $\{U\}$ 和 $\{VUV^\dagger\}$ 给出相同的最终链矩阵），因此答案显然是肯定的。对于库仑规范和朗道规范，答案是否定的——规范固定条件有多个解。这些就是著名的 Gribov 拷贝 [42]。关于这一主题近期的理论综述见 [43]。

固定到库仑规范或朗道规范的算法大多是局域的。局域算法对每个格点生成一个规范矩阵 $v_i(x)$，它使从该格点出发的链矩阵之和（由式~\ref{eq:Coulombgauge} 定义）最大化。然后将这一规范变换施加到格点上。重复这两个步骤直到找到整体极大值。有序乘积 $\prod_i v_i(x)$ 给出 $V(x)$。在实际应用中，我们将规范固定算法的这个终点称为 Gribov 拷贝。在大多数情况下，尤其是对于局域规范固定算法，它是 $\lVert A^V \rVert$ 的最近极值点，而非整体极小值。可以生成一组 Gribov 拷贝：先对组态施行一个不同的随机规范变换，然后应用相同的规范固定算法，或者使用不同的规范固定算法。关于不同算法有效性的近期研究见 [44,45]。

蒙特卡洛更新过程使得任何有限的统计样本中包含规范等价组态的可能性极低。因此不必担心规范拷贝。在规范固定的路径积分中，先验地人们希望每个组态只对一个 Gribov 拷贝求平均，并选择「最光滑」的那一个。找到给出最光滑组态的 $V(x)$ 已被证明是一个非平凡的问题。

悬而未决的问题是：（i）如何定义规范固定算法的「最光滑」解（即定义基本模域），（ii）如何确定一个解是唯一的（位于基本模域内部或其边界上），以及（iii）如何找到它。

最终的问题是——人们想要计算的量对 Gribov 拷贝的选择有怎样的依赖？已经进行了许多研究 [46]，但结果及其解释尚无定论。首先，尚不清楚能否对所有的组态以及在所研究的规范耦合范围内经济地找到位于基本模域内的最光滑解。在没有最光滑解的情况下，也不清楚（i）残余规范自由度是否仅仅对关联函数中的噪声有贡献，以及（ii）是否应对每个组态的一组 Gribov 拷贝求平均以减小这种噪声。在我看来，这些仍然是具有挑战性的问题，诸位中可能有人愿意去研究。

当人们想做格点微扰论，或通过传递矩阵形式体系取路径积分的哈密顿极限时，Gribov 拷贝和 Faddeev-Papov 构造（规范固定项和 Faddeev-Papov 行列式）的问题会再次出现。这一构造与连续情形中的构造类似 [36]。最后，让我提一下，在固定规范中，格点理论只有 BRST 不变性 [27]。虽然 BRST 不变性足以证明理论是可重整化的，但 BRST 不变性所允许的任一给定维数的算符集合，比具有完全规范不变性的集合更大。在作用量与算符改进的分析中，必须考虑的正是在这个扩大了的集合 [47]。
